\documentclass[12pt]{article}
\usepackage[utf8]{inputenc}
\usepackage[T1]{fontenc}
\usepackage{hyperref}
\title{Kapinaw\'a Language Resources for K-12 Curriculum Development}
\author{Compiled by Claude Sonnet 5.5 (claude-sonnet-5-5) for Andr\'as Kornai}
\date{}
\begin{document}
\maketitle

\section{Introduction}
Kapinaw\'a (ISO 639-3 xpn; Glottocode kapi1248) is the presumed ancestral language of the Kapinaw\'a people of Bu\'ique, Pernambuco, Brazil. It is unattested: no word list or text is known, and the community now speaks Portuguese. The Kapinaw\'a are said to descend from the Parati\'o, whose language is itself poorly attested, known only from two word lists published in 1958. Glottolog suggests it was related to Xukuru. A curriculum that uses the native language is therefore not possible on present evidence. What exists is Portuguese-language material on Kapinaw\'a culture, history and education.

\section{Primary Educational Resources}
No language primers or textbooks exist. The sources below concern Kapinaw\'a schooling and culture, in Portuguese.
\begin{itemize}
\item ``Caminhando pela cultura e lutas dos povos Kapinaw\'a em Bu\'ique-PE: um relato de campo'' (UFPE journal \emph{Mutir\~ao}). \url{https://periodicos.ufpe.br/revistas/mutiro/article/view/254320}
\item Same article, PDF download. \url{https://periodicos.ufpe.br/revistas/mutiro/article/download/254320/43354/220513}
\item A paper in \emph{Revista Anthropol\'ogicas} (UFPE) surfaced in the results. I did not verify its content. \url{https://periodicos.ufpe.br/revistas/revistaanthropologicas/article/download/23671/19327/46669}
\item A UFRPE final-year thesis (TCC) by Joelma Carla da Silva surfaced in the results. I did not verify its content. \url{https://repository.ufrpe.br/bitstream/123456789/5673/1/tcc_joelmacarladasilva.pdf}
\item A conference paper from anais.ueg.br surfaced in the results. I did not verify its content. \url{https://anais.ueg.br/index.php/spi/article/download/16952/13554/51374}
\end{itemize}

\section{Linguistic and Reference Materials}
\begin{itemize}
\item Wikipedia, ``Kapinaw\'a language'': states that the language is unattested and its classification unknown. \url{https://en.wikipedia.org/wiki/Kapinaw%C3%A1_language}
\item Wikipedia, ``Parati\'o language'': the probable ancestor, known from two 1958 word lists. \url{https://en.wikipedia.org/wiki/Parati%C3%B3_language}
\item Wikipedia, ``Kapinaw\'a'': the people. \url{https://en.wikipedia.org/wiki/Kapinaw%C3%A1}
\item BnF RAMEAU authority record. It notes that there are no longer any speakers of the language. \url{https://catalogue.bnf.fr/ark:/12148/cb18115807g}
\end{itemize}

\section{Religious and Community Resources}
No community organisation website or religious material was found in the results.

\section{Digital and Audio Resources}
None found. No recordings or dictionaries of Kapinaw\'a exist in the results.

\section{Curriculum Development Context}
The sources note that Indigenous rights in Brazil include teaching in mother tongues. In the Kapinaw\'a case the mother tongue is Portuguese. A school unit could cover Kapinaw\'a history, territory and struggle, drawing on the Portuguese-language scholarship above, with the community leading. Any language revitalization would require the community's own decision and new documentation, since no attested material exists to teach from.
\end{document}